% !PRINTED 492
\weberpage{492}

Then the products $\omega \omega_1, \omega \omega_2, \ldots, \omega \omega_m$ can, with the help of the equations (2), be put into the form
\[
\omega \omega_r = c_{r,1} \omega_1 + c_{r,2} \omega_2 + \cdots + c_{r,m} \omega_m \quad r = 1, 2, \ldots, m,
\]
in which the coefficients $c_{s,r}$ are integral rational numbers. But if one eliminates the $\omega_1, \omega_2, \ldots, \omega_m$ from these equations, it follows that
\[
\begin{vmatrix}
c_{1,1} - \omega, & c_{1,2}, & \ldots, & c_{1,m} \\
c_{2,1}, & c_{2,2} - \omega, & \ldots, & c_{2,m} \\
\cdots & \cdots & \cdots & \cdots \\
c_{m,1}, & c_{m,2}, & \ldots, & c_{m,m} - \omega
\end{vmatrix}
= 0,
\]
which, developed, contains the form
\[
\omega^m + C_1 \omega^{m-1} + \cdots + C_m = 0,
\]
in which the $C_1, C_2, \ldots, C_m$ are likewise integral rational numbers. But this shows that $\omega$ is an integer, as was to be proved.

3. \gesperrt{If $f(x)$ is a function irreducible in the field of rational numbers, and if a root $\Theta$ of $f(x) = 0$ is an integer, then all roots of $f(x)$ are integers.}

For if a rational function $F(x)$ vanishes for $x = \alpha$, then $F(x)$ is divisible by $f(x)$, and all roots of $f(x)$ are at the same time roots of $F(x)$ (Bd.\,I, \S.\,141). Now if $\alpha$ is an integer, there is a function
\[
F(x) = x^m + A_1 x^{m-1} + \cdots + A_m
\]
with integral coefficients $A_1, \ldots, A_m$ which vanishes for $x = \alpha$, and $F(x)$ therefore also vanishes for all the other roots of $f(x)$, which accordingly are all integers.

If
\[
f(x) = x^n + a_1 x^{n-1} + \cdots + a_n,
\]
then the $a_1, a_2, \ldots, a_n$ are composed by multiplication and addition from the roots of $f(x)$, and are therefore, by 2., integral rational numbers. From this it follows:

4. \gesperrt{If $\Theta$ is an integral algebraic number, then the equation of lowest degree $f(\Theta) = 0$ [\S.\,132, (2)] has integral coefficients.}
